\section{Algoritmos Implementados}
\label{sec:algoritmos}

\subsection{Detector de Falhas}

Nós que não são coordenadores executam um laço periódico. A cada intervalo,
enviam \texttt{Heartbeat} ao coordenador conhecido. Uma resposta válida zera o
contador de falhas consecutivas. Um timeout, uma resposta inválida ou um
coordenador ausente incrementa esse contador. Ao atingir o limiar configurado,
o nó inicia uma eleição e zera o contador.

\subsection{Eleição Bully}

O iniciador envia \texttt{StartElection} aos nós de identificador maior. Um
destino maior responde com \texttt{accepted=true} e inicia sua própria eleição
quando necessário. Se nenhum nó maior responder, o iniciador atualiza seu
coordenador local e envia \texttt{AnnounceCoordinator} aos demais nós.

O acesso ao estado compartilhado é protegido por um lock. Assim, a atualização
do identificador da eleição, do coordenador e do indicador de eleição ativa é
atômica em relação aos callbacks RPC e ao laço de heartbeat. O tratamento de
eleições antigas impede que anúncios atrasados substituam uma eleição mais
recente.

\subsection{Pseudocódigo}

O comportamento do detector pode ser resumido como:
\par\medskip
\par\smallskip
\begin{lstlisting}[language={}]
enquanto o processo estiver ativo:
    se este no nao for o coordenador:
        enviar Heartbeat com deadline
        se resposta valida: falhas <- 0
        caso contrario: falhas <- falhas + 1
        se falhas >= limiar:
            iniciar_eleicao()
            falhas <- 0
    aguardar intervalo
\end{lstlisting}
